% !TEX program = pdflatex
\documentclass[12pt]{article}

\usepackage{sbc-template}
\usepackage{graphicx,url}
\usepackage[utf8]{inputenc}
\usepackage[brazilian]{babel}
\usepackage{amsmath}

\sloppy

\title{Implementação e Comparação de Algoritmos de Ordenação}

\author{Antônio Gemesson Sousa\inst{1}, Gleidson Luan Sena\inst{1}, Ivan Vitor Dias\inst{1}}

\address{Departamento de Computação -- Universidade Federal do Piauí (UFPI)\\
  Teresina -- PI -- Brasil
  \email{\{autor1, autor2, ivan.oliveira\}@ufpi.edu.br}
}

\begin{document}

\maketitle

\begin{abstract}
  % Escrever em inglês, máximo 10 linhas.
  This paper presents the implementation and experimental comparison of four
  sorting algorithms: Insertion Sort, Merge Sort, Counting Sort and Radix Sort.
  The algorithms are evaluated over random integer arrays of varying sizes and
  value ranges, relating experimental results to their theoretical complexity
  analyses.
\end{abstract}

\begin{resumo}
  % Máximo 10 linhas.
  Este artigo apresenta a implementação e a comparação experimental de quatro
  algoritmos de ordenação: Insertion Sort, Merge Sort, Counting Sort e Radix
  Sort. Os algoritmos são avaliados sobre vetores de inteiros aleatórios com
  diferentes tamanhos e intervalos de valores, relacionando os resultados
  experimentais às respectivas análises de complexidade teórica.
\end{resumo}

% ============================================================
\section{Introdução}
% ============================================================

% - Contextualizar ordenação e sua importância.
% - Diferenciar algoritmos de comparação e lineares.
% - Apresentar o objetivo do trabalho.

\section{Descrição e Análise Teórica}
% ============================================================

\subsection{Insertion Sort}

O Insertion Sort é um algoritmo de ordenação por comparação que constrói a
sequência ordenada de forma incremental, inserindo cada novo elemento na posição
correta em relação aos elementos já ordenados. O algoritmo mantém um subvetor
ordenado à esquerda e, a cada iteração, retira o próximo elemento não ordenado
e o desloca para a posição correta dentro desse subvetor, movendo os elementos
maiores uma posição à frente.

\textbf{Análise de complexidade:}
\begin{itemize}
  \item \textbf{Melhor caso -- $O(n)$:} Ocorre quando o vetor já está ordenado.
    Em cada iteração, o elemento corrente é maior ou igual ao último elemento do
    subvetor ordenado, e nenhuma troca é realizada. O algoritmo executa apenas
    as $n - 1$ comparações do laço externo.
  \item \textbf{Caso médio -- $O(n^2)$:} Para entradas com distribuição
    aleatória, espera-se que, em média, cada elemento precise ser deslocado
    metade do subvetor ordenado à sua esquerda. O número de comparações e
    deslocamentos cresce quadraticamente com $n$.
  \item \textbf{Pior caso -- $O(n^2)$:} Ocorre quando o vetor está ordenado em
    ordem decrescente. Cada elemento inserido precisa ser comparado e deslocado
    em relação a todos os elementos já ordenados, resultando em $O(n^2)$
    operações.
\end{itemize}

\subsection{Merge Sort}

O Merge Sort é um algoritmo de ordenação por comparação baseado na estratégia
de divisão e conquista. O vetor é dividido recursivamente ao meio até que cada
subvetor contenha um único elemento (trivialmente ordenado). Em seguida, os
subvetores são mesclados dois a dois em ordem crescente, produzindo subvetores
cada vez maiores até que o vetor completo esteja ordenado. A etapa de mesclagem
percorre os dois subvetores simultaneamente, copiando sempre o menor elemento
para o vetor resultante.

\textbf{Análise de complexidade:}
\begin{itemize}
  \item \textbf{Melhor, médio e pior caso -- $O(n \log n)$:} Independentemente
    da distribuição inicial dos dados, o Merge Sort sempre divide o vetor em
    duas metades e mescla os resultados. A árvore de recursão possui altura
    $O(\log n)$ e, em cada nível, o trabalho total de mesclagem é $O(n)$.
    Portanto, a complexidade é $O(n \log n)$ em todos os casos.
  \item \textbf{Espaço auxiliar -- $O(n)$:} A etapa de mesclagem requer vetores
    auxiliares para armazenar temporariamente os elementos, resultando em uso
    de memória proporcional ao tamanho da entrada.
\end{itemize}

\subsection{Counting Sort}

O Counting Sort é um algoritmo de ordenação linear que não realiza comparações
entre elementos. Ele opera contando o número de ocorrências de cada valor
distinto no vetor de entrada e, a partir de uma soma de prefixo sobre o vetor
de contagem, determina a posição exata de cada elemento no vetor de saída. A
iteração do vetor de saída de trás para frente garante a estabilidade do
algoritmo.

\textbf{Análise de complexidade:}
\begin{itemize}
  \item \textbf{Todos os casos -- $O(n + k)$:} onde $n$ é o número de elementos
    e $k$ é o valor do maior elemento (o intervalo dos valores). O algoritmo
    realiza uma passada de $O(n)$ para encontrar o máximo, uma de $O(n)$ para
    contar as ocorrências, uma de $O(k)$ para calcular a soma de prefixo e uma
    de $O(n)$ para reconstruir o vetor de saída.
\end{itemize}

\textbf{Parâmetros além do tamanho do vetor:}
Ao contrário dos algoritmos de comparação, a complexidade do Counting Sort
depende não apenas de $n$, mas também de $k$, o valor máximo presente no vetor.
Quando $k = O(n)$, o algoritmo é efetivamente linear. Porém, quando $k$ cresce
muito além de $n$ -- como nos intervalos $[0, 100n]$ ou $[0, 1000n]$ previstos
nos experimentos --, o custo de $O(k)$ passa a dominar, degradando o desempenho.
Além disso, o Counting Sort requer $O(k)$ de memória auxiliar, o que pode ser
proibitivo para valores de $k$ muito elevados.

\subsection{Radix Sort}

O Radix Sort é um algoritmo de ordenação linear que ordena os elementos dígito
a dígito, do menos significativo ao mais significativo. Em cada passada, os
elementos são distribuídos com base no dígito corrente usando uma versão estável
do Counting Sort (com soma de prefixo), preservando a ordem relativa de
elementos com o mesmo dígito. A estabilidade é essencial para que as passadas
anteriores não sejam desfeitas. O dígito na posição \texttt{exp} é extraído por
\texttt{(arr[i] / exp) \% 10}.

\textbf{Análise de complexidade:}
\begin{itemize}
  \item \textbf{Todos os casos -- $O(d \cdot n)$:} onde $d$ é o número de
    dígitos do maior elemento e $n$ é o número de elementos. Cada uma das $d$
    passadas executa em $O(n)$, pois o Counting Sort sobre dígitos de 0 a 9
    tem $k = 10$ fixo. Como $d = O(\log k)$, a complexidade total é
    $O(n \log k)$.
  \item \textbf{Vantagem sobre o Counting Sort:} O Radix Sort é menos sensível
    à magnitude dos valores, pois o custo por passada depende apenas de $n$ e
    não de $k$ diretamente.
\end{itemize}

% ============================================================
\section{Metodologia}
% ============================================================

% Preencher com as informações do ambiente de testes:
Os experimentos foram conduzidos nas seguintes condições:

\begin{itemize}
  \item \textbf{Linguagem de programação e compilador:} C++, % compilador e versão
  \item \textbf{Sistema operacional:} % ex: Ubuntu 24.04
  \item \textbf{Processador:} % ex: Intel Core i5-1135G7, 2.4 GHz
  \item \textbf{Memória RAM:} % ex: 8 GB DDR4 3200 MHz
\end{itemize}

Para cada algoritmo, foram utilizados vetores de inteiros gerados aleatoriamente
com distribuição uniforme. Foram testados os seguintes tamanhos de entrada $n$:
% ex: 1.000, 10.000, 100.000, 500.000 e 1.000.000

Para cada tamanho $n$, os experimentos consideraram quatro intervalos de valores:
$[0, n]$, $[0, 10n]$, $[0, 100n]$ e $[0, 1000n]$. Para cada combinação de $n$
e intervalo, foram gerados 5 vetores distintos, e cada algoritmo foi executado
5 vezes por vetor. O tempo de execução reportado corresponde à média das
execuções. Os algoritmos foram avaliados sequencialmente, sem paralelismo.

% ============================================================
\section{Resultados e Discussão}
% ============================================================

% Inserir tabelas e/ou gráficos com os resultados.
% Discutir:
% - diferenças entre algoritmos de mesma categoria
% - diferenças entre categorias (comparação vs. linear)
% - influência do tamanho n e do intervalo de valores
% - relação com a análise teórica
% - quando os lineares têm vantagem ou desvantagem

% ============================================================
\section{Conclusão}
% ============================================================

% Resumir os principais achados e relacionar com os objetivos do trabalho.

\bibliographystyle{sbc}
\bibliography{sbc-template}

\end{document}
